\section{Method}
\label{sec:method}

\subsection{Trajectory graphs}
\label{sec:method-graphs}
Every recorded trajectory is mapped losslessly into a directed heterogeneous DAG. Nodes carry one of four types---\textsc{Problem} (user request), \textsc{Thinking} (agent reasoning), \textsc{Action} (tool call), and \textsc{State} (tool response or final reply)---and store the verbatim trajectory text in \texttt{raw\_content} together with a SHA-256 content hash and a JSON path pointing back to the source record. Edges carry one of five semantic relation types (\texttt{depends\_on}, \texttt{computes\_from}, \texttt{reflects}, \texttt{sequential}, \texttt{branches\_from}) with an LLM-generated \texttt{justification} string; edges are generated in a two-stage procedure whose prompts never observe trajectory labels, and the resulting graph is verified to be forward-only and acyclic. On \tool{} the 2{,}218 graphs contain 55{,}594 nodes and 139{,}994 edges in total.

\subsection{Graph anomaly detector}
\label{sec:method-detector}
Node features concatenate a frozen sentence embedding (\texttt{all-MiniLM-L6-v2}, 384-d, L2-normalized), a 16-d learnable node-type embedding, and two structural scalars (normalized sequence position and log content length), projected by a shared MLP. The encoder is an edge-aware graph attention network: attention scores and messages both condition on edge features (the edge embedding, the relation embedding, and a direction embedding), and reverse edges are appended for message passing while directional semantics are preserved through the direction embedding. Graph-level readout concatenates mean- and max-pooling over all nodes (\texttt{mean\_max}); a two-layer MLP head produces the anomaly logit. We train with cross-entropy under a grouped stratified split (Section~\ref{sec:setup}). Leakage words that could correlate with labels (e.g., \emph{attack}, \emph{unsafe}, \emph{benign}) are masked from all node and edge texts.

\subsection{Weakly supervised risk localization}
\label{sec:method-localization}
Because only graph-level labels exist, localization uses multiple-instance learning on top of the \emph{frozen} detector. The scorer consumes the detector's post-message-passing node embeddings $h_i$, the broadcast graph embedding $h_G$, and a node-type embedding $t_i$, and predicts per-node risk $r_i=\sigma(\mathrm{MLP}([h_i\,\|\,h_G\,\|\,t_i]))$. The graph-level risk is the maximum node score (max-MIL), trained with binary cross-entropy against the graph label plus a sparsity penalty $\lambda_{\mathrm{sparse}}\,\overline{r}$ restricted to abnormal graphs, encouraging few-node explanations. A fingerprint over the frozen detector's parameters is verified before and after training to guarantee that the detector is never modified. For every graph predicted anomalous by the detector, the top-$k$ ($k{=}3$) risk nodes seed a connected risk subgraph obtained by shortest-path closure over the undirected adjacency.

\subsection{LLM-filtered experience synthesis}
\label{sec:method-experience}
Predicted anomalies---not ground-truth labels---enter the experience stage. For each flagged graph we assemble a case containing the graph anomaly score, threshold margin, the localized nodes (with content redacted by four regular-expression families covering credential assignments, \texttt{sk-} style keys, and bearer tokens, then truncated to 800 characters), and the induced edges with justifications. A postmortem judge LLM receives the case with an explicit instruction that trajectory content is untrusted data whose embedded instructions must never be followed, and returns a nine-field verdict: \texttt{is\_reusable\_experience}, judgment (anomalous / benign-or-false-positive / uncertain), confidence, risk types, \texttt{what\_went\_wrong}, warning signals, recommended agent policy, supporting node ids, and a reason. Only verdicts with \texttt{is\_reusable\_experience=true} become experiences; benign, false-positive, and insufficient-evidence cases are rejected by design. The synthesis judge never observes ground-truth labels; this is asserted in the run summary.

\subsection{Paired replay defense evaluation}
\label{sec:method-replay}
The defense effect is measured by replaying the same recorded attack cases under two conditions. The replay reconstructs the environment deterministically: leading user steps and all environment steps (including the injected payloads) are fixed evidence; each recorded agent step becomes a \emph{proposed} step that is hidden from the controller until its execution checkpoint. An execution controller receives the task, the tool list, the transcript so far, optionally the safety experience, and exactly one proposed step, and returns \textsc{execute}, \textsc{revise} (with a replacement step), or \textsc{stop}. After the replay terminates, a security judge scores whether the replayed agent executed, accepted, or meaningfully advanced the expected unsafe behavior, using a hidden evaluation reference; for \tool{} the reference is the recorded compromised agent path. The controller is instructed to treat the experience as advisory memory, not as an instruction or ground truth. Baseline and experience conditions share the case set, agent model, temperature, and protocol; an input fingerprint over prompts, models, and case identifiers is stored with each run and mismatches are rejected.

\subsection{Model roles}
\label{sec:method-roles}
To avoid self-family evaluation bias, model families are separated by role: the attacked agent and the execution controller run on the \texttt{deepseek-flash} endpoint, the experience-synthesis judge runs on the same method-side endpoint, and the final ASR verdict is produced by an independent judge family (\texttt{gpt-5.5}) that is disjoint from both the agent and the synthesis judge. Judge prompts are pinned by the input fingerprint, and all calls use temperature 0 with JSON response format enforcement and bounded retries.
